% !TEX program = pdflatex
\documentclass[11pt]{article}
\usepackage[margin=1in]{geometry}
\usepackage{amsmath,amssymb,amsthm,mathtools,bm}
\usepackage{physics}
\usepackage{hyperref}
\usepackage{cleveref}
\usepackage{xcolor}
\usepackage{listings}
\usepackage{microtype}
\lstdefinestyle{py}{
  language=Python,
  basicstyle=\ttfamily\footnotesize,
  numbers=left,
  numberstyle=\tiny, 
  stepnumber=1,
  numbersep=6pt,
  frame=single,
  breaklines=true,
  showstringspaces=false,
  tabsize=2,
  columns=fullflexible
}
\lstset{style=py}

\title{Neural Nexus: A Minimal, Well-Posed Framework for Noisy Variational Quantum Algorithms}
\author{Ryan O. Van Gelder \\ \\ Citizen Gardens}
\date{\today}

\begin{document}
\maketitle

\begin{abstract}
We present a concise, mathematically coherent specification for a noisy variational quantum algorithm (VQA) framework suitable for near-term devices. The formulation defines Hilbert spaces, density-operator dynamics under completely positive trace-preserving (CPTP) maps, task losses, gradient estimation via parameter-shift, and a minimal two-qubit demonstration with reproducible baselines. We include reference implementations and validation protocols.
\end{abstract}

\section{Preliminaries}
Let \(n\in\mathbb N\) be the number of qubits and \(\mathcal H=(\mathbb C^{2})^{\otimes n}\). States are density operators \(\rho\in\mathcal D(\mathcal H)=\{\rho\succeq0,\ \mathrm{Tr}\,\rho=1\}\). A quantum channel is a CPTP map \(\mathcal E: \mathcal D(\mathcal H)\to\mathcal D(\mathcal H)\) with Kraus form \(\mathcal E(\rho)=\sum_k E_k\rho E_k^{\dagger}\), \(\sum_k E_k^{\dagger}E_k=I\) \cite{NielsenChuang,BreuerPetruccione}.

Controls are real parameters \(\theta\in\mathbb R^{p}\). We use a hardware-efficient unitary ansatz \(U(\theta)\in\mathrm{U}(2^n)\) compiled from native rotations and entanglers \cite{CerezoVQAReview}.

\section{Open dynamics and noise}
We evolve for \(K\) steps of size \(\Delta t\). Each step applies the unitary followed by noise:
\begin{equation}
  \rho_{k+1}=\mathcal N_{\Delta t}\!\big( U(\theta)\,\rho_k\,U(\theta)^{\dagger}\big),\qquad k=0,\dots,K-1.
\end{equation}
Per-qubit noise composes amplitude damping and dephasing \cite{NielsenChuang}:
\begin{align}
  \mathcal A_{T_1}(\rho) &= K_0\rho K_0^{\dagger}+K_1\rho K_1^{\dagger},\quad K_0=\ket{0}\!\bra{0}+\sqrt{1-\eta}\,\ket{1}\!\bra{1},\ K_1=\sqrt{\eta}\,\ket{0}\!\bra{1},\ \eta=1-e^{-\Delta t/T_1},\\
  \mathcal D_{\phi}(\rho) &= L_0\rho L_0^{\dagger}+L_1\rho L_1^{\dagger},\quad L_0=\sqrt{1-p}\,I,\ L_1=\sqrt{p}\,Z,\ p=\tfrac{1-e^{-\Delta t/T_\phi}}{2}.
\end{align}
The two-qubit channel is the tensor product of single-qubit Kraus operators. A small correlated component can be added using a two-qubit depolarizing channel \(\mathcal E_{\mathrm{dep}}(\rho)=(1-\varepsilon)\rho+\varepsilon\,I/4\).

\section{Tasks}
\paragraph{Noisy state preparation.} Given a pure target \(\ket{\psi_*}\), define \(\rho_T\) after \(K\) steps and loss
\(L(\theta)=1-\langle\psi_*|\rho_T|\psi_*\rangle\).
\paragraph{Coherence preservation.} For observables \(\mathcal O\), use 
\(L(\theta)=\frac{1}{|\mathcal O|}\sum_{O\in\mathcal O}\big|\mathrm{Tr}(O\rho_T)-\mathrm{Tr}(O\rho_0)\big|\).

\section{Ansatz}
We use a depth-\(d\) hardware-efficient circuit with local \(R_Z\) and \(R_X\) rotations per qubit followed by entangling CNOTs on a fixed graph.

\section{Gradients}
For gates generated by Pauli operators, the parameter-shift rule yields exact gradients using two shifted evaluations per parameter \cite{MariPS,SchuldPS}. For noisy evolution with parameter-independent channels, the rule applies by placing the channel after the shifted gate. Stochastic SPSA is an alternative when shifts are not available \cite{SpallSPSA}.

\section{Two-qubit demonstration}
We target \(\ket{\Phi^+}=(\ket{00}+\ket{11})/\sqrt{2}\) starting from \(\ket{00}\!\bra{00}\). Fidelity is estimated via Pauli expectations
\begin{equation}
F=\tfrac{1+\langle XX\rangle-\langle YY\rangle+\langle ZZ\rangle}{4},
\end{equation}
or via a projector measurement. Classical shadows provide an efficient estimator for larger systems \cite{HuangShadows}.

\section{Validation protocol}
We require: (i) parameter-shift vs finite-difference agreement within tolerance, (ii) baselines (noiseless, identity, random) reported, (iii) sweeps over \(T_1,T_\phi,K\) with fidelity vs iterations and shots, and (iv) CPTP and unitarity sanity checks.

\section{Resources}
Per iteration with parameter-shift and \(p\) parameters, shot usage is \(m(1+2p)\). Gate counts per loss evaluation scale with depth and steps as described in the spec.

\section{Reference implementation (selected snippets)}
\subsection*{Pauli helpers and ansatz}
\begin{lstlisting}
import numpy as np
I = np.array([[1,0],[0,1]], complex)
X = np.array([[0,1],[1,0]], complex)
Y = np.array([[0,-1j],[1j,0]], complex)
Z = np.array([[1,0],[0,-1]], complex)
PM = {'I': I, 'X': X, 'Y': Y, 'Z': Z}

def pauli_matrix(label: str):
    a, b = label
    return np.kron(PM[a], PM[b])

def rz(theta):
    return np.array([[np.exp(-1j*theta/2), 0], [0, np.exp(1j*theta/2)]], complex)

def rx(theta):
    return np.array([[np.cos(theta/2), -1j*np.sin(theta/2)],
                     [-1j*np.sin(theta/2), np.cos(theta/2)]], complex)

def cx():
    return np.array([[1,0,0,0],[0,1,0,0],[0,0,0,1],[0,0,1,0]], complex)

def hardware_efficient_ansatz(theta, n_qubits=2, depth=2):
    theta = theta.reshape(depth, n_qubits, 2)
    U = np.eye(4, complex); CNOT = cx()
    for layer in range(depth):
        U1 = np.kron(rz(theta[layer,0,0]) @ rx(theta[layer,0,1]),
                     rz(theta[layer,1,0]) @ rx(theta[layer,1,1]))
        U = CNOT @ U1 @ U
    return U
\end{lstlisting}

\subsection*{CPTP local noise and correlated depolarizing}
\begin{lstlisting}
def _single_qubit_kraus(T1, T_phi, dt):
    eta = 1 - np.exp(-dt / T1)
    K0 = np.array([[1, 0],[0, np.sqrt(1-eta)]], complex)
    K1 = np.array([[0, np.sqrt(eta)],[0, 0]], complex)
    p = (1 - np.exp(-dt / T_phi)) / 2
    L0 = np.sqrt(1 - p) * I
    L1 = np.sqrt(p) * Z
    return [L @ K for L in (L0,L1) for K in (K0,K1)]

def local_noise(rho, T1, T_phi, dt):
    K1 = _single_qubit_kraus(T1, T_phi, dt)
    out = np.zeros_like(rho)
    for A in K1:
        for B in K1:
            K = np.kron(A, B)
            out += K @ rho @ K.conj().T
    return out

def two_qubit_depolarizing(rho, epsilon):
    dim = rho.shape[0]
    return (1 - epsilon) * rho + (epsilon / dim) * np.eye(dim, complex)
\end{lstlisting}

\subsection*{Loss, gradients, and optimization}
\begin{lstlisting}
RHO00 = np.diag([1,0,0,0]).astype(complex)
phi = np.array([1/np.sqrt(2), 0, 0, 1/np.sqrt(2)], complex)
PHIp = np.outer(phi, phi.conj())

def loss_function(theta, shots=None, T1=50.0, T_phi=25.0, dt=0.1, K=4, noise_fn=None):
    U = hardware_efficient_ansatz(theta, 2, 2)
    rho = RHO00.copy()
    for _ in range(K):
        rho = U @ rho @ U.conj().T
        rho = (noise_fn or (lambda r: local_noise(r, T1, T_phi, dt)))(rho)
    F = float(np.real(np.trace(rho @ PHIp)))
    if shots is not None:
        F = np.clip(F + np.random.normal(0, 1/np.sqrt(shots)), 0, 1)
    return 1 - F

def ps_gradient(theta, L):
    g = np.zeros_like(theta)
    for j in range(len(theta)):
        e = np.zeros_like(theta); e[j] = 1.0
        g[j] = 0.5*(L(theta + np.pi/2*e) - L(theta - np.pi/2*e))
    return g

def optimize_vqa(noise_fn, theta0, max_iter=100, eta=0.1, **loss_kwargs):
    theta = theta0.copy(); history = []
    L = lambda th: loss_function(th, noise_fn=noise_fn, **loss_kwargs)
    for _ in range(max_iter):
        history.append(L(theta))
        theta -= eta * ps_gradient(theta, L)
    return theta, history
\end{lstlisting}

\subsection*{Parameter-shift validation and demo}
\begin{lstlisting}
def validate_parameter_shift(theta, j, loss_fn, rtol=1e-2, atol=1e-3, shots=None):
    test_loss = (lambda th: loss_fn(th, shots=shots)) if shots else loss_fn
    shift = np.pi/2
    e = np.zeros_like(theta); e[j] = 1.0
    grad_ps = 0.5*(test_loss(theta + shift*e) - test_loss(theta - shift*e))
    h = 1e-4
    grad_fd = (test_loss(theta + h*e) - test_loss(theta - h*e)) / (2*h)
    if not np.isclose(grad_ps, grad_fd, rtol=rtol, atol=atol):
        raise ValueError(f"parameter-shift check failed: ps={grad_ps}, fd={grad_fd}")
    return grad_ps

def run_2qubit_demo():
    theta0 = np.random.normal(0, 0.1, size=8)
    T1, T_phi, dt, K = 50.0, 25.0, 0.1, 4
    noise_variants = {
        'local_only': lambda rho: local_noise(rho, T1, T_phi, dt),
        'with_correlated': lambda rho: two_qubit_depolarizing(local_noise(rho, T1, T_phi, dt), 0.05),
    }
    results = {}
    for name, nf in noise_variants.items():
        Lps = lambda th, shots=None: loss_function(th, shots=shots, T1=T1, T_phi=T_phi, dt=dt, K=K, noise_fn=nf)
        validate_parameter_shift(theta0, 0, Lps, shots=10000)
        theta_opt, hist = optimize_vqa(nf, theta0, max_iter=50, eta=0.1, T1=T1, T_phi=T_phi, dt=dt, K=K)
        F = 1 - loss_function(theta_opt, T1=T1, T_phi=T_phi, dt=dt, K=K, noise_fn=nf)
        results[name] = F
    return results
\end{lstlisting}

\section{Limitations}
The current demo is restricted to two qubits and tensor-factorized noise with a simple correlated depolarizing component. Extending to larger systems requires efficient state representations and measurement allocation strategies.

\section{Conclusion}
The framework is executable, verifiable, and scalable in a controlled manner. It provides a fast path to proof on real devices while maintaining mathematical rigor.

\begin{thebibliography}{9}
\bibitem{NielsenChuang}
M. A. Nielsen and I. L. Chuang, \emph{Quantum Computation and Quantum Information}, 10th Anniversary Edition, Cambridge University Press, 2010.

\bibitem{BreuerPetruccione}
H.-P. Breuer and F. Petruccione, \emph{The Theory of Open Quantum Systems}, Oxford University Press, 2002.

\bibitem{CerezoVQAReview}
M. Cerezo et al., ``Variational Quantum Algorithms,'' \emph{Nature Reviews Physics} 3, 625–644 (2021).

\bibitem{MariPS}
A. Mari, T. R. Bromley, J. Izaac, M. Schuld, and N. Killoran, ``Estimating the gradient and higher-order derivatives on quantum hardware,'' \emph{PRX Quantum} 2, 030301 (2021).

\bibitem{SchuldPS}
M. Schuld, S. Sweke, and J. J. Meyer, ``Effect of data encoding on the expressive power of quantum models,'' \emph{Phys. Rev. A} 103, 032430 (2021). [Contains parameter-shift discussion.]

\bibitem{SpallSPSA}
J. C. Spall, ``Multivariate stochastic approximation using a simultaneous perturbation gradient approximation,'' \emph{IEEE Trans. Automatic Control} 37(3), 332–341 (1992).

\bibitem{HuangShadows}
H.-Y. Huang, R. Kueng, and J. Preskill, ``Predicting many properties of a quantum state from few measurements,'' \emph{Nature Physics} 16, 1050–1057 (2020).
\end{thebibliography}

\end{document}